\section{Text, lists and math}

% Frame title and optional subtitle: \begin{frame}{Title}{Subtitle}.
% \section{...} starts a section: its name appears at the top right of the
% slides after it and in the Overview.
\begin{frame}{Lists}{itemize, enumerate and description}
  \begin{columns}[T]
    \begin{column}{0.48\textwidth}
      \begin{itemize}
        \item A bullet point
        \item Another one, with \textbf{bold}, \emph{emphasis} and
              \textcolor{bannerblue}{colour}
          \begin{itemize}
            \item A nested point
          \end{itemize}
      \end{itemize}
      \begin{enumerate}
        \item A numbered point
        \item Another numbered point
      \end{enumerate}
    \end{column}
    \begin{column}{0.48\textwidth}
      \begin{description}
        \item[Term] Its explanation
        \item[Other] Another explanation
      \end{description}
      {\small Size commands work too: \cmd{small}, \cmd{footnotesize}, \ldots}
    \end{column}
  \end{columns}
\end{frame}

% Inline math \( ... \) or $ ... $, display math \[ ... \], and the AMS
% environments (align, equation, ...). The class option (11pt in main.tex)
% scales all text, math included.
\begin{frame}{Math and blocks}
  Inline math such as \( Re = \rho U D / \mu \) sits in the text; displayed
  equations get their own line:
  \[ \rho \left( \frac{\partial \mathbf{u}}{\partial t}
       + \mathbf{u} \cdot \nabla \mathbf{u} \right)
     = -\nabla p + \mu \nabla^2 \mathbf{u} \]
  \begin{align}
    m \ddot{y} + c \dot{y} + k y &= F_y(t) \\
    \omega_n &= \sqrt{k / m}
  \end{align}
  \begin{block}{A block}
    Blocks frame a key message: \cmd{begin\{block\}\{Title\}} \ldots
  \end{block}
\end{frame}
